{\small
\setlength{\tabcolsep}{3.0pt}
\renewcommand{\arraystretch}{1.18}
\begin{tabular}{L{4.6cm}L{9.0cm}}
\toprule
Элемент проверки & Как выполнено \\
\midrule
Целевая переменная & $y$ — готовность рекомендовать чаще используемый сервис другим студентам, шкала 1--10. \\
Матрица данных & Использованы 60 строк объединённого листа ответов; признаки взяты из блока характеристик экспертов и очищенной таблицы \texttt{cleaned\_survey\_with\_scores.csv}. \\
Подбор коэффициентов & Перед оцениванием признаки стандартизированы; целевая переменная оставлена в исходной шкале 1--10. Коэффициенты найдены методом наименьших квадратов. \\
Проверка модели & Оценены $R^2$, скорректированный $R^2$, F-критерий общей значимости, t-критерии для коэффициентов, p-value, VIF и график «факт--прогноз». \\
Интерпретация & Модель используется как объяснительная, а не как прогнозная; учитывались устойчивые связи с p-value < 0{,}05. \\
\bottomrule
\end{tabular}}